% alt: a chain of eighteen sites with an interval T of three sites in the middle, a wider interval R of seven sites around it, the two end pieces of the chain beyond the next widening labelled C, two smaller end pieces inside them labelled E, and a brace under the gap between E and T marked D_p log L.
% Section 8.3 (The one-dimensional separation lemma), Lemma lem:logsep and its proof: R = T^{+(w+1)}, C = Lambda \ T^{+(2w+2)}, R_0 = T, C_0 = E with dist(T,E) >= D_p log L >= 3w+3, so that the reset lemma applies with I(R:C) <= 2S(R) <= b (drawn for w = 1, so R adds two sites and C starts four sites beyond T on each side); the dots are the sites and the upward legs the physical indices C^q, every leg is open (nothing is contracted); the red enclosure is T (three sites), the blue enclosure is R = T^{+(w+1)} (seven sites), the two dashed enclosures are the two pieces of C = Lambda \ T^{+(2w+2)} (three sites at the left end, four at the right), the two inner enclosures at the ends are the two pieces of E, a subset of C, and the brace under the five sites between E and T marks the separation dist(T,E) >= D_p log L.
\begin{tenkz}[lattice={1x18}, frame=plane, physical=up, bonds=none]
  \tn{} & \tn{} & \tn{} & \tn{} & \tn{} & \tn{} & \tn{} & \tn{} & \tn{} & \tn{} & \tn{} & \tn{} & \tn{} & \tn{} & \tn{} & \tn{} & \tn{} & \tn{}
  \tnmark[form=enclosure, species=selected, label pos=s]{(1,8) .. (1,10)}{$T$}
  \tnmark[form=enclosure, species=secondary, label pos=n]{(1,6) .. (1,12)}{$R=T^{+(w+1)}$}
  \tnmark[form=enclosure, species=complement, label pos=n]{(1,1) .. (1,3)}{$C$}
  \tnmark[form=enclosure, species=complement, label pos=n]{(1,15) .. (1,18)}{$C$}
  \tnmark[form=enclosure, species=collar, label pos=s]{(1,1) .. (1,2)}{$E$}
  \tnmark[form=enclosure, species=collar, label pos=s]{(1,17) .. (1,18)}{$E$}
  \tnmark[form=bracket, species=collar, label pos=s]{(1,3) .. (1,7)}{$D_p\log L$}
\end{tenkz}
